\chapter{Stress-Energy Tensor OPE and Virasoro Algebra}

\section{Stress-Energy OPE, Central Charge and Ward Identities}

As we have established, the holomorphic component of the stress tensor, \(T(z)\), generates infinitesimal conformal transformations and controls the local Ward identities. The operator product expansion of \(T(z)\) with itself, \(T(w)\), plays a distinguished role in two-dimensional conformal field theory. It determines:
\begin{itemize}
    \item the algebra of infinitesimal conformal transformations,
    \item the quantum modification of the classical Witt algebra,
    \item the appearance of the central charge.
\end{itemize}
This is precisely the point where the quantum nature of the theory manifests itself most clearly.

\subsection{General Form of the \(T(z)T(w)\) OPE}

The stress-energy tensor \(T_{\mu\nu}\) is a conserved current. Under rescaling, its scale dimension and spin are fixed to be:
\[
    \begin{aligned}
         & \Delta_T = 2,         &  & s_T = 2,          \\
         & \Delta_{\bar{T}} = 2, &  & s_{\bar{T}} = -2,
    \end{aligned}
\]
which immediately imply the conformal weights:
\[
    (h_T, \bar{h}_T) = (2,0) \;, \quad (h_{\bar{T}}, \bar{h}_{\bar{T}}) = (0,2).
\]
As a quasi-primary field, it rescales globally as \(T(\lambda z) = \lambda^2 T(z)\). This fixes its two-point correlation function to be strictly of the form:
\[
    \langle T(z)T(w) \rangle = \frac{\text{const.}}{(z-w)^4}.
\]
By convention, this constant is denoted as \(c/2\), where \(c\) is called the \textbf{central charge}.

The full OPE of \(T(z)T(w)\) for \(z \to w\) is tightly constrained. Both scale invariance and the requirement that correlation functions of local fields must be meromorphic\footnote{Meromorphic means that the function is analytic except for isolated poles.} dictate the general structure of the singular terms:
\[
    T(z)T(w) = \frac{c/2}{(z-w)^4} + \frac{J(w)}{(z-w)^3} + \frac{Q(w)}{(z-w)^2} + \frac{R(w)}{z-w} + \text{regular terms},
\]
where the generic operators must have specific holomorphic dimensions: \(h_J = 1\), \(h_Q = 2\), and \(h_R = 3\), since the total conformal weight of the singular terms must match the sum of the weights of the two \(T\) operators, which is \(4\).

Because \(T(z)\) is a bosonic field, the OPE must be symmetric under the exchange of coordinates and operators, meaning \(T(z)T(w) = T(w)T(z)\). We can write the reversed OPE for \(w \to z\):
\[
    T(w)T(z) = \frac{c/2}{(w-z)^4} + \frac{J(z)}{(w-z)^3} + \frac{Q(z)}{(w-z)^2} + \frac{R(z)}{w-z} + \text{regular terms}.
\]
Assuming analyticity of the fields for \(z\) near \(w\), we can express the right-hand side entirely at the coordinate \(w\) by Taylor expanding the fields \(J(z)\) and \(Q(z)\) around \(w\). For instance:
\[
    J(z) = J(w) + (z-w)\partial_w J(w) + \frac{1}{2}(z-w)^2 \partial_w^2 J(w) + \dots
\]
Substituting these expansions into the reversed OPE, paying attention to the sign changes from the denominators, we get:
\[
    T(w)T(z) = \frac{c/2}{(z-w)^4} - \frac{J(w)}{(z-w)^3} + \frac{Q(w) + \partial_w J(w)}{(z-w)^2} + \frac{\partial_w Q(w) + \frac{1}{2}\partial_w^2 J(w) - R(w)}{z-w} + \text{regular terms}.
\]
By matching the singularities term by term with the original \(T(z)T(w)\) expansion (since the OPE is symmetric), we find strict algebraic constraints on the unknown fields:
\[
    J(w) \equiv 0 \;, \quad R(w) = \frac{1}{2}\partial_w Q(w),
\]
which means that \(J(w)\) along with its derivatives must vanish, and \(R(w)\) is not an independent operator but is determined by \(Q(w)\).

If we now make the highly reasonable physical hypothesis that \uline{\(T(z)\) is the \textit{only} conformal field of spin and dimension 2 in the algebra of local fields},\footnote{This is a very reasonable assumption, since if another conformal field \(Q(z)\) had \(h=2\) such that its 2-point correlation function enforced \(\langle Q(w)Q(z)\rangle = \frac{c^{\prime} / 2}{(z-w)^4}\) then we could add it to \(T\) to create from \(T+Q\) a larger algebra with central charge \(c + c'\).} then \(Q(w)\) must be proportional to \(T(w)\) itself. Indeed, we know that \(h_Q=2\) and it must be holomorphic thus chiral:
\[
    \partial_{\bar{w}} Q(w) = 0, \quad \partial_w \bar{Q}(\bar{w}) = 0,
\]
which looks exactly like the structure of \(T(w)\). We can thus write \(Q(w) = \kappa T(w)\). The OPE then simplifies to:
\[
    T(z)T(w) = \frac{c/2}{(z-w)^4} + \frac{\kappa T(w)}{(z-w)^2} + \frac{\frac{\kappa}{2} \partial_w T(w)}{z-w} + \text{regular terms}.
\]

To fix the proportionality constant \(\kappa\), we must look at the physical meaning of the second-order pole in an OPE with the stress-energy tensor.
\begin{itemize}
    \item For any general field, the coefficient of the \((z-w)^{-2}\) pole strictly corresponds to its conformal weight \(h\), which dictates its transformation under the global dilatation generator \(L_0\).
    \item Since we analytically know that the stress-energy tensor \(T(w)\) is defined as a field of conformal dimension \(h = 2\), requiring dimensional consistency under a global dilatation \(z, w \to \lambda z, \lambda w\) immediately fixes \(\kappa = 2\).
    \item Furthermore, we know the OPE of the stress-energy tensor with a generic primary field \(\Phi(w)\) of weight \(h\) must contain a second-order pole with coefficient \(h\Phi(w)\):
          \[
              T(z)\Phi(w) = \frac{h\Phi(w)}{(z-w)^2} + \frac{\partial_w \Phi(w)}{z-w} + \text{regular terms},
          \]
          which is the defining OPE for a primary field. Now, \(T(w)\) is a quasi-primary field of weight \(h = 2\), so intuitively, the second-order pole in the \(T(z)T(w)\) OPE must also have coefficient \(2T(w)\) to be consistent with the transformation properties of \(T\) itself.
\end{itemize}

This consequently sets the coefficient of the first-order pole to \(1\), yielding the fundamental \(TT\) operator product expansion valid for any two-dimensional Conformal Field Theory:
\begin{equation}
    T(z)T(w) = \frac{c/2}{(z-w)^4} + \frac{2T(w)}{(z-w)^2} + \frac{\partial_w T(w)}{z-w} + \text{regular terms}.
    \label{eq:TT_OPE_CFT2}
\end{equation}
This is the general form of the \(T(z)T(w)\) OPE, which holds for all two-dimensional CFTs. It encodes the algebraic structure of the conformal symmetry and the presence of the central charge \(c\), which characterizes the quantum anomaly in the conformal symmetry.

\subsubsection{Central Charge}

The constant \(c\) is a fundamental characteristic of each specific theory, rightfully named the \textbf{central charge}. It is fixed once the stress-energy tensor is defined, which in turn is known when the underlying Lagrangian is specified.

As we have seen in our practical calculations for free models, different theories possess different central charges: for a free massless boson \(c = 1\), while for a free massless Majorana fermion \(c = 1/2\).

For physical, unitary models defined by Hermitian Hamiltonians, the Hermiticity of the stress-energy tensor \(T\) requires that its two-point correlator be strictly positive, leading to the bound:
\[
    c > 0 \quad \text{for Hermitian (unitary) models}.
\]
An entirely analogous \(\bar{T}\bar{T}\) OPE defines an anti-holomorphic central charge \(\bar{c}\). If the theory is parity invariant (meaning it is completely symmetric under the spatial reflection \(z \leftrightarrow \bar{z}\)), the two central charges must coincide:
\[
    c = \bar{c}.
\]

\subsubsection{Variation of \(T\) under Conformal Transformations}

To understand how the energy-momentum tensor transforms, we first recall the Cauchy integral formula and its generalization for derivatives:
\[
    f(z) = \oint_z \frac{\mathrm{d} w}{2\pi i} \frac{f(w)}{w - z} \;, \quad \frac{\mathrm{d}^n f(z)}{\mathrm{d}z^n} = n! \oint_z \frac{\mathrm{d} w}{2\pi i} \frac{f(w)}{(w - z)^{n+1}},
\]
where the contour integral is taken over a path in the complex \(w\)-plane encircling the point \(z\).

We also recall that the infinitesimal variation of any local conformal field \(A(z, \bar{z})\) under a generic holomorphic conformal transformation \(\epsilon(z)\) is generated by its commutator with the conserved charge (from section \ref{sec:noether_OPE_classification}), which can be expressed as a contour integral of the OPE:
\[
    \delta_\epsilon A(z, \bar{z}) = \oint_z \frac{\mathrm{d} w}{2\pi i} \epsilon(w) T(w) A(z, \bar{z}).
\]
We now apply this general formula to the specific case where the field itself is the stress-energy tensor, \(A(z, \bar{z}) = T(z)\). By inserting the fundamental \(TT\) operator product expansion (expanded around \(z\)) into the integral:
\[
    T(w)T(z) \sim \frac{c/2}{(w - z)^4} + \frac{2T(z)}{(w - z)^2} + \frac{\partial_z T(z)}{w - z},
\]
we obtain:
\[
    \delta_\epsilon T(z) = \oint_z \frac{\mathrm{d} w}{2\pi i} \epsilon(w) \left[ \frac{c/2}{(w - z)^4} + \frac{2T(z)}{(w - z)^2} + \frac{\partial_z T(z)}{w - z} \right].
\]
Using the Cauchy derivative formulas on the test function \(\epsilon(w)\), we can evaluate each term of the singular expansion exactly. The simple pole yields \(\epsilon(z) \partial_z T(z)\), the double pole yields \(2 (\partial_z \epsilon(z)) T(z)\), and the fourth-order pole produces \(\frac{1}{3!} \partial_z^3 \epsilon(z)\) multiplied by \(c/2\). Putting it all together, we find:
\[
    \delta_\epsilon T(z) = \epsilon(z)\partial_z T(z) + 2(\partial_z \epsilon(z))T(z) + \frac{c}{12}\partial_z^3 \epsilon(z).
\]
This represents the infinitesimal conformal transformation of \(T(z)\). The first two terms perfectly match the expected transformation law for a primary field of conformal weight \(h = 2\). However, the third term is an inhomogeneous anomaly proportional to the central charge \(c\), explicitly demonstrating that the stress-energy tensor is only a quasi-primary field, not a primary one.

By an exponentiation procedure which we will not report here, one can integrate this infinitesimal variation to obtain the exact finite transformation law under a macroscopic conformal map \(z \mapsto w(z)\):
\[
    T(z) = \left(\frac{\partial w}{\partial z}\right)^2 T(w) + \frac{c}{12} \{w, z\}_S,
\]
where the symbol \(\{w, z\}_S\) denotes the \textbf{Schwarzian derivative}, defined as:
\[
    \{w, z\}_S = \frac{\partial_z^3 w}{\partial_z w} - \frac{3}{2} \left( \frac{\partial_z^2 w}{\partial_z w} \right)^2.
\]

\begin{example}[Infinitesimal Limit of the Schwarzian Transformation]
    We want to expand the finite transformation law for a small coordinate change \(w = z + \epsilon(z)\) to reproduce the infinitesimal variation \(\delta_\epsilon T(z)\).

    First, we compute the derivatives of \(w\) with respect to \(z\):
    \(\partial_z w = 1 + \partial_z \epsilon\), \(\partial_z^2 w = \partial_z^2 \epsilon\), and \(\partial_z^3 w = \partial_z^3 \epsilon\).

    For the Jacobian prefactor, we expand to first order in \(\epsilon\):
    \[
        \left(\frac{\partial w}{\partial z}\right)^2 = (1 + \partial_z \epsilon)^2 \approx 1 + 2\partial_z \epsilon.
    \]
    For the Schwarzian derivative, the denominator is \(1 + \partial_z \epsilon \approx 1\), so to first order:
    \[
        \{w, z\}_S \approx \frac{\partial_z^3 \epsilon}{1} - \frac{3}{2} (\partial_z^2 \epsilon)^2 \approx \partial_z^3 \epsilon.
    \]
    (The squared term \((\partial_z^2 \epsilon)^2\) is second-order in \(\epsilon\) and thus vanishes).

    Now, substituting these into the finite transformation law \(T(z) = (\partial_z w)^2 T(w) + \frac{c}{12} \{w, z\}_S\):
    \[
        T(z) \approx (1 + 2\partial_z \epsilon) T(z + \epsilon) + \frac{c}{12} \partial_z^3 \epsilon.
    \]
    Taylor expanding \(T(z + \epsilon) \approx T(z) + \epsilon \partial_z T(z)\) and keeping only first-order terms:
    \[
        T(z) \approx (1 + 2\partial_z \epsilon) (T(z) + \epsilon \partial_z T(z)) + \frac{c}{12} \partial_z^3 \epsilon,
    \]
    \[
        T(z) \approx T(z) + \epsilon \partial_z T(z) + 2(\partial_z \epsilon)T(z) + \frac{c}{12} \partial_z^3 \epsilon.
    \]
    Subtracting \(T(z)\) from both sides isolates the variation \(\delta_\epsilon T(z)\):
    \[
        \delta_\epsilon T(z) = \epsilon(z) \partial_z T(z) + 2(\partial_z \epsilon(z)) T(z) + \frac{c}{12} \partial_z^3 \epsilon(z).
    \]
    This exactly matches the result derived from the contour integral of the \(TT\) OPE.
\end{example}

\subsection{Physical Meaning of the Central Charge}

To fully grasp the physical significance of the central charge \(c\), we investigate how the vacuum of the theory behaves when we move from the infinite flat plane to a space with a finite geometry. We do this by mapping the complex \(z\)-plane to a cylinder of circumference \(L\) using the standard conformal transformation:
\[
    z = e^{\frac{2\pi}{L}w} \quad \text{or equivalently} \quad w = \frac{L}{2\pi} \log z,
\]
where \(w = \tau + i\sigma\) is the coordinate on the cylinder.

To see how the stress-energy tensor transforms under this map, we first compute the necessary derivatives: \(\partial_z w = \frac{L}{2\pi z}\), \(\partial_z^2 w = \frac{-L}{2 \pi z^2}\) and finally \(\partial_z^3 w = \frac{L}{\pi z^3}\). The Schwarzian derivative \(\{w, z\}_S\) measures the ``non-Möbius'' nature of the transformation, and by plugging the derivatives of the logarithm into its definition, the scale factor \(L/2\pi\) neatly cancels out, leaving the following expression for the cylinder:
\[
    \{w, z\}_S = \frac{1}{2z^2}.
\]
Substituting these ingredients into the finite transformation law for the stress-energy tensor, we obtain the exact form of the tensor on the cylinder:
\[
    T_{\text{cyl}}(w) = \left(\frac{2\pi}{L}\right)^2 \left\{ T_{\text{pl}}(z)z^2 - \frac{c}{24} \right\}.
\]
We can now evaluate the vacuum expectation value of this operator. On the infinite plane, global conformal invariance dictates that the one-point function of any quasi-primary field with non-zero dimension must vanish. Therefore, \(\langle T_{\text{pl}}(z) \rangle = 0\).

However, inserting this vanishing value into our transformation yields a strictly non-zero expectation value on the cylinder:
\[
    \langle T_{\text{cyl}}(w) \rangle = -\frac{\pi^2 c}{6L^2}.
\]
This is a profound physical result. The expectation value of the diagonal component \(T + \bar{T} = T_{11}\) is directly proportional to the vacuum energy density. While the vacuum energy is exactly zero on the infinite plane, the periodic boundary conditions of the cylindrical geometry induce a macroscopic shift in the zero-point energy.

This non-zero vacuum energy density is the conformal field theory realization of the \textbf{Casimir energy}, caused entirely by finite-size boundary effects. Including both the holomorphic and anti-holomorphic sectors, the total Casimir energy density is given by:
\[
    \langle T_{11} \rangle = \frac{1}{2\pi} \{ \langle T \rangle + \langle \bar{T} \rangle \} = -\frac{\pi c}{6L^2}.
\]
This explicitly demonstrates that the central charge \(c\) is not just an abstract algebraic parameter governing the Virasoro anomaly, but a measurable physical quantity that dictates the finite-size scaling behavior of the theory's ground state.

\section{Ward Identities and Correlation Functions}

We now derive the conformal Ward identities. Consider a correlation function of local fields:
\[
    \langle \Phi_1(z_1, \bar{z}_1) \cdots \Phi_n(z_n, \bar{z}_n) \rangle.
\]
Inserting the stress tensor \(T(z)\) into the correlator, and using its role as the generator of conformal transformations, one finds the local Ward identity:
\[
    \langle T(z) \Phi_1(z_1, \bar{z}_1) \cdots \Phi_n(z_n, \bar{z}_n) \rangle = \sum_{i=1}^n \left[ \frac{h_i}{(z - z_i)^2} + \frac{1}{z - z_i} \partial_{z_i} \right] \langle \Phi_1(z_1, \bar{z}_1) \cdots \Phi_n(z_n, \bar{z}_n) \rangle,
\]
provided that all the fields \(\Phi_i\) are primary fields.
Similarly, in the anti-holomorphic sector:
\[
    \langle \bar{T}(\bar{z}) \Phi_1(z_1, \bar{z}_1) \cdots \Phi_n(z_n, \bar{z}_n) \rangle = \sum_{i=1}^n \left[ \frac{\bar{h}_i}{(\bar{z} - \bar{z}_i)^2} + \frac{1}{\bar{z} - \bar{z}_i} \partial_{\bar{z}_i} \right] \langle \Phi_1(z_1, \bar{z}_1) \cdots \Phi_n(z_n, \bar{z}_n) \rangle.
\]
These fundamental relations are known as the local conformal Ward identities.

\begin{proof}
    The derivation follows straightforwardly from the Operator Product Expansion (OPE). Since the stress-energy tensor \(T(z)\) is purely holomorphic everywhere away from operator insertions, the correlator \(\langle T(z) \Phi_1 \cdots \Phi_n \rangle\) behaves as a meromorphic function of the coordinate \(z\). Its singularities arise solely when \(z\) approaches one of the insertion points \(z_i\). By applying the defining OPE between the stress tensor and a primary field for each \(\Phi_i\), and subsequently summing all the singular contributions, one exactly recovers the expression for the local Ward identity.
\end{proof}

\subsubsection{Global Ward Identities}

If we restrict our attention to the three globally defined holomorphic conformal transformations (the Möbius transformations), which correspond to the choices \(\epsilon(z) = 1\), \(\epsilon(z) = z\), and \(\epsilon(z) = z^2\), the local Ward identities reduce to a set of global constraints:
\[
    \sum_{i=1}^n \partial_{z_i} \langle \Phi_1(z_1, \bar{z}_1) \cdots \Phi_n(z_n, \bar{z}_n) \rangle = 0,
\]
\[
    \sum_{i=1}^n (z_i \partial_{z_i} + h_i) \langle \Phi_1(z_1, \bar{z}_1) \cdots \Phi_n(z_n, \bar{z}_n) \rangle = 0,
\]
\[
    \sum_{i=1}^n (z_i^2 \partial_{z_i} + 2 h_i z_i) \langle \Phi_1(z_1, \bar{z}_1) \cdots \Phi_n(z_n, \bar{z}_n) \rangle = 0.
\]
There are perfectly analogous equations governing the anti-holomorphic sector.

These relations constitute the Ward identities for the global conformal group in two dimensions. They impose stringent differential constraints on the correlation functions of primary fields, strictly reflecting how these fields must transform under global conformal mappings. Physically, each identity enforces the invariance of the theory under a specific subgroup: the first equation guarantees translation invariance, the second enforces invariance under dilations and rotations, and the third corresponds to special conformal transformations. Solving this set of differential equations is a crucial step in understanding the structure of two-dimensional Conformal Field Theories, as it allows us to completely determine the exact analytical form of two- and three-point correlation functions up to overall multiplicative constants.

\subsection{Correlation Functions from Ward Identities}

Let us derive the exact spatial dependence of the two-point function of primary fields using the global conformal Ward identities. Consider the general two-point correlation function:
\[
    G(z_1, \bar{z}_1; z_2, \bar{z}_2) = \langle \Phi_1(z_1, \bar{z}_1) \Phi_2(z_2, \bar{z}_2) \rangle.
\]
We systematically apply the constraints imposed by the unbroken global conformal subgroup (translations, dilatations, and special conformal transformations) to restrict the form of this function.

\paragraph{1. Translational Invariance:}
Invariance under rigid spatial translations implies that the correlation function can only depend on the relative coordinate differences:
\[
    z_{12} = z_1 - z_2, \quad \bar{z}_{12} = \bar{z}_1 - \bar{z}_2.
\]
Thus, the function reduces to \( G(z_1, \bar{z}_1; z_2, \bar{z}_2) = G(z_{12}, \bar{z}_{12}) \).

\paragraph{2. Dilatation Invariance:}
Invariance under global scale transformations imposes a first-order differential equation on the correlator. From the dilatation Ward identity, we obtain:
\[
    (z_{12}\partial_{z_{12}} + h_1 + h_2)G(z_{12}, \bar{z}_{12}) = 0.
\]
The general solution to this differential equation fixes the holomorphic coordinate dependence up to an arbitrary function of the anti-holomorphic coordinate:
\[
    G(z_{12}, \bar{z}_{12}) = \frac{f(\bar{z}_{12})}{z_{12}^{h_1 + h_2}}.
\]
Applying the exact same logic to the anti-holomorphic sector tightly restricts \( f(\bar{z}_{12}) \), leaving us with a single undetermined overall constant \( C_{12} \):
\[
    G(z_{12}, \bar{z}_{12}) = \frac{C_{12}}{z_{12}^{h_1 + h_2} \bar{z}_{12}^{\bar{h}_1 + \bar{h}_2}}.
\]

\paragraph{3. Special Conformal Invariance:}
Finally, the Ward identity for special conformal transformations imposes an incredibly restrictive algebraic condition. It dictates that the constant \( C_{12} \) must vanish identically unless the conformal weights of the two fields are strictly equal:
\[
    h_1 = h_2, \quad \text{and} \quad \bar{h}_1 = \bar{h}_2.
\]
By absorbing this mathematical requirement into Kronecker deltas, we arrive at the final, exact form of the two-point function for primary fields:
\[
    \langle \Phi_1(z_1, \bar{z}_1) \Phi_2(z_2, \bar{z}_2) \rangle = \frac{C_{12} \delta_{h_1, h_2} \delta_{\bar{h}_1, \bar{h}_2}}{z_{12}^{2h_1} \bar{z}_{12}^{2\bar{h}_1}}.
\]
If the fields are purely scalar (\( h = \bar{h} = \Delta/2 \)), this beautifully simplifies to the familiar physical form \( C_{12}/|z_{12}|^{2\Delta} \). This proves that two-point functions are non-zero \uline{only} if the fields share the exact same scaling dimension and conformal spin.

\subsubsection{Connection to the Stress-Energy Tensor OPE}

As we established in earlier chapters, these global geometric constraints are actually a direct, macroscopic consequence of the fundamental Operator Product Expansion (OPE) between the stress-energy tensor \( T(z) \) and the primary fields.

Inserting \( T(z) \) into a general \( n \)-point correlator yields the fundamental Conformal Ward Identity:
\[
    \langle T(z) \Phi_1(z_1, \bar{z}_1) \cdots \Phi_n(z_n, \bar{z}_n) \rangle = \sum_{i=1}^n \left[ \frac{h_i}{(z-z_i)^2} + \frac{1}{z-z_i} \partial_{z_i} \right] \langle \Phi_1(z_1, \bar{z}_1) \cdots \Phi_n(z_n, \bar{z}_n) \rangle.
\]
The differential operators generating translations (\( \partial_z \)), dilatations (\( z\partial_z + h \)), and special conformal transformations natively emerge by expanding this identity and integrating over closed contours surrounding the field insertions. This elegantly unifies the local algebraic structure of the Virasoro modes with the macroscopic macroscopic symmetries of the correlation functions we just derived.

\subsubsection{Three-Point Function from Ward Identities}

A similar, albeit slightly more tedious, application of the exact same global Ward identities completely fixes the coordinate dependence of the three-point function. For three generic primary fields, conformal invariance strictly requires the highly symmetric form:
\[
    \langle \Phi_1(z_1, \bar{z}_1) \Phi_2(z_2, \bar{z}_2) \Phi_3(z_3, \bar{z}_3) \rangle = \frac{C_{123}}{z_{12}^{h_1 + h_2 - h_3} z_{23}^{h_2 + h_3 - h_1} z_{31}^{h_3 + h_1 - h_2} \bar{z}_{12}^{\bar{h}_1 + \bar{h}_2 - \bar{h}_3} \bar{z}_{23}^{\bar{h}_2 + \bar{h}_3 - \bar{h}_1} \bar{z}_{31}^{\bar{h}_3 + \bar{h}_1 - \bar{h}_2}},
\]
where we define \( z_{ij} = z_i - z_j \) and \( \bar{z}_{ij} = \bar{z}_i - \bar{z}_j \).

Notice that the entire spatial dependence is perfectly and exclusively determined by the conformal weights of the operators. The only remaining unknown in this equation is the overall multiplicative factor \( C_{123} \), which is known as the \textbf{structure constant} (or OPE coefficient).

Because global conformal symmetry strictly fixes all two- and three-point functions up to these constants, \uline{solving a conformal field theory formally reduces to finding two discrete sets of data}:
\begin{enumerate}
    \item The \textbf{spectrum} of the theory: identifying all allowed primary fields and their conformal weights \( (h_i, \bar{h}_i) \).
    \item The \textbf{structure constants} \( C_{ijk} \): calculating the numbers that govern the exact interaction strengths in the Operator Product Expansion.
\end{enumerate}
Once the spectrum and the structure constants are known, any higher-order correlation function (like the four-point functions) can, in principle, be systematically reduced to these fundamental building blocks by recursively applying the OPE.